% ============================================================
%  Chapter 3 — Mathematical Model
% ============================================================
\section{Mathematical Model}
\label{ch:model}

This chapter develops the \MILP{} model for operational speed
optimisation under \CII{} constraints. \Cref{sec:model:problem} describes
the problem. \Cref{sec:model:notation} introduces the notation.
\Cref{sec:model:base} presents the base formulation.
\Cref{sec:model:flexible} describes the flexible schedule extension.
\Cref{sec:model:fuel} specifies the fuel consumption model.
\Cref{sec:model:cii} details the \CII{} constraint derivation.

% ── 3.1 Problem description ───────────────────────────────────
\subsection{Problem Description}
\label{sec:model:problem}

A single \RoRo{} vessel operates a cyclic service visiting a fixed set of
ports. The vessel makes one port call at each port per cycle and must
depart each port within a published time window. The operator selects a
sailing speed for each voyage leg (arc) from a finite discrete set. The
objective is to minimise total voyage cost --- comprising fuel cost and,
where applicable, \ETS{} carbon cost --- subject to \CII{} compliance,
schedule adherence, and physical speed limits.

The model is a \MILP{} because speed is represented as a discrete
variable (\ie{} an integer choice from a finite speed grid), while the
\CII{} constraint introduces a ratio structure that is linearised by
appropriate reformulation.

% ── 3.2 Notation ──────────────────────────────────────────────
\subsection{Notation}
\label{sec:model:notation}

\subsubsection{Sets}
\label{sec:model:sets}

\begin{itemize}
  \item $\setN$: set of port nodes.
  \item $\setA \subseteq \setN \times \setN$: set of directed arcs
    $(i,j)$ representing voyage legs.
  \item $\setV$: finite set of discrete speed levels (knots),
    $\setV = \{v_{\min}, \ldots, v_{\max}\}$.
\end{itemize}

\subsubsection{Parameters}
\label{sec:model:params}

\begin{itemize}
  \item $\dist$: sailing distance on arc $(i,j)$, in nautical miles.
  \item $\cargo$: cargo (in tonnes) carried on arc $(i,j)$; fixed
    exogenously at the mean observed load.
  \item $\fuelrate$: fuel consumption on arc $(i,j)$ at speed $v$,
    in kg/nm; pre-computed from the vessel's polynomial fuel curve
    (see \cref{sec:model:fuel}).
  \item $\budget$: maximum allowed transit time on arc $(i,j)$,
    in hours; derived from the published port schedule.
  \item $C_{\mathrm{fuel}}$: fuel price in \euro{}/tonne.
  \item $C_{\mathrm{CO}_2}$: \ETS{} carbon price in \euro{}/kg\,CO$_2$.
  \item $\gamma$: CO$_2$ emission factor (kg\,CO$_2$ per tonne of fuel).
  \item $\ciithr$: \CII{} threshold in \gGTnm{}; scenario-specific
    (see \cref{ch:scenarios}).
  \item $\GT{}$: vessel gross tonnage.
  \item $f_i, f_m$: \IMO{} correction factors for ice class and
    capacity, respectively.
\end{itemize}

\subsubsection{Decision Variables}
\label{sec:model:vars}

\begin{equation}
  \xvar \in \{0,1\} \quad \forall (i,j) \in \setA,\ v \in \setV
  \label{eq:defx}
\end{equation}

$\xvar$ equals 1 if the vessel sails arc $(i,j)$ at speed $v$, and 0
otherwise. Each arc must be sailed at exactly one speed.

% ── 3.3 Base formulation ──────────────────────────────────────
\subsection{Base Formulation (Fixed Schedule)}
\label{sec:model:base}

\subsubsection{Objective Function}
\label{sec:model:obj}

The objective minimises total voyage cost comprising fuel cost and
\ETS{} carbon cost:

\begin{equation}
  \min \quad
  \sum_{(i,j)\in\setA} \sum_{v\in\setV}
    \xvar \cdot \dist \cdot \fuelrate
    \cdot \left( \frac{C_{\mathrm{fuel}}}{1000}
                 + \frac{C_{\mathrm{CO}_2} \cdot \gamma}{1000} \right)
  \label{eq:obj}
\end{equation}

\noindent where the factor $1/1000$ converts kg to tonnes for fuel cost
and accounts for units in the carbon cost term.

\subsubsection{Constraints}
\label{sec:model:constraints}

\paragraph{Speed assignment.}
Each arc is sailed at exactly one speed:

\begin{equation}
  \sum_{v \in \setV} \xvar = 1 \quad \forall (i,j) \in \setA
  \label{eq:speedassign}
\end{equation}

\paragraph{Schedule (time budget) constraints.}
The transit time on each arc must not exceed the published time window:

\begin{equation}
  \sum_{v \in \setV} \xvar \cdot \frac{\dist}{v} \leq \budget
  \quad \forall (i,j) \in \setA
  \label{eq:schedule}
\end{equation}

\paragraph{CII constraint.}
The attained \CII{} must not exceed the scenario threshold
$\ciithr$:

\begin{equation}
  \frac{\displaystyle\sum_{(i,j)\in\setA} \sum_{v\in\setV}
          \xvar \cdot \dist \cdot \fuelrate \cdot \gamma}
       {f_i \cdot f_m \cdot \GT \cdot
        \displaystyle\sum_{(i,j)\in\setA} \sum_{v\in\setV}
          \xvar \cdot \dist}
  \leq \ciithr
  \label{eq:cii_frac}
\end{equation}

Because the denominator of \cref{eq:cii_frac} depends on the decision
variables, the constraint is non-linear as written. However, since
$\sum_{v} \xvar = 1$ implies $\sum_{v} \xvar \cdot \dist = \dist$ for
each arc, the total distance $D = \sum_{(i,j)} \dist$ is in fact
\textit{constant} (independent of speed choices). The constraint
therefore reduces to a linear form:

\begin{equation}
  \sum_{(i,j)\in\setA} \sum_{v\in\setV}
    \xvar \cdot \dist \cdot \fuelrate \cdot \gamma
  \leq \ciithr \cdot f_i \cdot f_m \cdot \GT \cdot D
  \label{eq:cii_linear}
\end{equation}

\noindent \cref{eq:cii_linear} is a single linear inequality, keeping the
overall model a \MILP{}.

\paragraph{Integrality.}
\begin{equation}
  \xvar \in \{0,1\} \quad \forall (i,j) \in \setA,\ v \in \setV
  \label{eq:integrality}
\end{equation}

\subsubsection{Complete Formulation}
\label{sec:model:complete}

The complete base model is:

\begin{align}
  \min \quad & \sum_{(i,j)\in\setA} \sum_{v\in\setV}
    \xvar \cdot \dist \cdot \fuelrate
    \cdot \left( \frac{C_{\mathrm{fuel}} + C_{\mathrm{CO}_2}\cdot\gamma}
                      {1000} \right)
    \tag{\ref{eq:obj}} \\
  \text{s.t.} \quad
  & \sum_{v\in\setV} \xvar = 1
    & \forall (i,j)\in\setA
    \tag{\ref{eq:speedassign}} \\
  & \sum_{v\in\setV} \xvar \cdot \frac{\dist}{v} \leq \budget
    & \forall (i,j)\in\setA
    \tag{\ref{eq:schedule}} \\
  & \sum_{(i,j)\in\setA} \sum_{v\in\setV}
      \xvar \cdot \dist \cdot \fuelrate \cdot \gamma
    \leq \ciithr \cdot f_i \cdot f_m \cdot \GT \cdot D
    \tag{\ref{eq:cii_linear}} \\
  & \xvar \in \{0,1\}
    & \forall (i,j)\in\setA,\ v\in\setV
    \tag{\ref{eq:integrality}}
\end{align}

% ── 3.4 Flexible schedule extension ───────────────────────────
\subsection{Flexible Schedule Extension}
\label{sec:model:flexible}

In the fixed-schedule formulation (\cref{sec:model:base}), per-arc time
budgets $\budget$ enforce the published timetable. In the flexible
schedule variant (Scenario~S5), the per-arc budget constraints
\eqref{eq:schedule} are replaced by a single \textit{global cycle-time
cap}:

\begin{equation}
  \sum_{(i,j)\in\setA} \sum_{v\in\setV} \xvar \cdot \frac{\dist}{v}
  \leq T_{\max}
  \label{eq:cycle_cap}
\end{equation}

\noindent where $T_{\max}$ is the maximum allowable total transit time
per complete loop. This relaxation allows the vessel to redistribute
sailing time freely across legs, enabling further slow-steaming on
legs that previously had tight time budgets.

% ── 3.5 Fuel consumption model ────────────────────────────────
\subsection{Fuel Consumption Model}
\label{sec:model:fuel}

The fuel consumption rate $\fuelrate$ (kg/nm) on arc $(i,j)$ at speed
$v$ is pre-computed from a vessel-specific sixth-degree polynomial
calibrated to three loading conditions:

\begin{equation}
  f^{(c)}(v) = \sum_{k=0}^{6} a_k^{(c)} \cdot v^k,
  \quad c \in \{\text{light},\ \text{medium},\ \text{full}\}
  \label{eq:poly_fuel}
\end{equation}

\noindent where the coefficients $a_k^{(c)}$ are derived from SCA's
Bunker Consumption data. For a given cargo load $\cargo$, the fuel rate
is linearly interpolated between adjacent loading conditions. The
polynomials are valid over the operational speed range
$[v_{\min}, v_{\max}] = [10, 16]$ knots.

\begin{remark}
  The sixth-degree polynomial exhibits a non-monotonic relationship
  between 12 and 13~knots for all three loading conditions. While the
  overall trend is increasing (higher speed $\Rightarrow$ more fuel),
  the polynomial dips slightly at 12~kn relative to 13~kn. This is
  a known artefact of polynomial fitting to empirical data and is
  acknowledged as a modelling limitation
  (see \cref{sec:disc:limitations}).
\end{remark}

% ── 3.6 CII constraint derivation ─────────────────────────────
\subsection{CII Constraint Derivation for RoRo Vessels}
\label{sec:model:cii}

This section documents the correct derivation of the \CII{} constraint
for a \RoRo{} vessel, which differs from formulations in the literature
that use deadweight tonnage or transport work.

\subsubsection{Emission Computation}

Total CO$_2$ emissions on a voyage are:

\begin{equation}
  E = \sum_{(i,j)\in\setA} \sum_{v\in\setV}
        \xvar \cdot \dist \cdot \fuelrate \cdot \gamma
  \quad [\text{kg CO}_2]
  \label{eq:total_co2}
\end{equation}

\noindent where $\gamma = 3{,}114$ kg\,CO$_2$/t\,fuel for heavy fuel
oil (HFO), following \IMO{} guidelines \parencite{mepc336}.

\subsubsection{CII Denominator}

For a \RoRo{} vessel, the \CII{} denominator is:

\begin{equation}
  \Omega = f_i \cdot f_m \cdot \GT \cdot D
  \label{eq:cii_denom}
\end{equation}

\noindent For M/V~Obbola: $\GT = 19{,}918$, $f_i = 1.0$,
$f_m = 1.05$ (ice class 1A correction), giving
$\Omega / D = 20{,}913.9$ GT per nautical mile.

\subsubsection{Attained CII}

\begin{equation}
  \ciiatt = \frac{E \times 1000}{\Omega}
  \quad [\text{g CO}_2 / (\GT \cdot \text{nm})]
  \label{eq:cii_attained}
\end{equation}

\noindent The factor of 1000 converts kg to grams.

\subsubsection{Relationship to Transport Work Formulations}

Common formulations in the literature define the \CII{} denominator as
deadweight tonnage times nautical miles, \ie{}
$\Omega_{\mathrm{DWT}} = \mathrm{DWT} \times D$. For \RoRo{} vessels,
this is incorrect under \IMO{} regulation. Using
$\mathrm{DWT} = 11{,}446$~t for M/V~Obbola in place of
$f_i \cdot f_m \cdot \GT = 20{,}914$ would overestimate the attained
\CII{} by a factor of $20{,}914 / 11{,}446 \approx 1.83$, leading to
an incorrect regulatory assessment. This distinction is a key finding
of this thesis and has been corrected relative to earlier model
versions.
